\documentclass{article}
\usepackage[english,russian]{babel}
\usepackage{mathtools}
\usepackage{listings}
\begin{document}
\author{Гоголев Иван}
\title{Критерии для оценки тестов}
\maketitle
\section{Введение}
    Для оценки различных подходов к генерации тестов, необходимо разработать критерии. Примеры рассчетов критериев
    для различных наборов тестов будем проводить на следующих примерах.

\begin{tabular}{|c|c|c|}
\hline
№ & Условие & Ограничения \\
\hline
1 & Найти сумму двух чисел a и b & \( 0 \leq a,b \leq 10^9\) \\
\hline
2 & Для данного простого p найти первообразный корень & \( 2 \leq p \leq 10^6\), p - простое \\
\hline
3 &  &  \\
\hline
\end{tabular}

    В соответствии с классичскими подходами к разработке тестов критерии должны отражать следующие свойства:
\begin{enumerate}
\item Наличие тестов, находящихся на границе пространства тестов
\item Ранномерность заполнения пространства тестов
\end{enumerate}

Назовем критерием функцию \(K: T \mapsto M, M \subseteq R \), где T - множество наборов тестов, R - множество
действительных чисел. Конкретный тестовый набор будем обозначать буквой t.



\section{Экспертная оценка}
    При ручной подготовке тестов основным критерием качества является оценка автора набора тестов \(K_1\):
    \( T \mapsto \{ 0, 1 \}  \). При этом работа по подготовке тестов продолжается пока \(K_1(t) \neq 1\)

    При оценке таким способом автоматической генерации тестов можно использовать более сложные функции
\begin{enumerate}
\item \(K_2 : T \mapsto \{ 0,..., n \}  \)
\item \(K_3 : T \mapsto [0, 1]   \)
\end{enumerate}
\section{Критерии неиспользующие семантику}
    Данные критерии используют для тестов модель черного ящика.
\subsection{Минимальный/максимальный размер теста}
    Пусть тестовый набор записан в памяти компьютера в некоторой кодировке тогда 
    \(K_4(t) := \min(S) \), где S - множество размеров тестов в памяти компьютера,
     аналогично определяется \(K_5\) - максимальный размер теста.
\subsection{Размах}
    Определим размах: \(K_6(t) := K_5(t) - K_4(t)\)
\subsection{Равномерность}
    Из равномерности распределения кодировок на множестве бинарных строк
    следует равномерность заполнения тестами пространства тестов. Простейшей проверкой этого является следующее:
    Пусть \(s_1 \leq s_2 \ldots \leq s_n \) - размеры тестов в отсортированном порядке. 
    Определим \(K_7 = \min(s_{i + 1} - s_i), i < n\).
\subsection{Пример рассчета}
    Рассмотрим рассчет критериев на нескольких наборов для задачи №2:
\begin{enumerate}
\item 2, 3, 5, 7
\item 2, 999863
\item 2, 999983, 99991, 99989, 9973, 997, 97, 31, 19, 7
\end{enumerate}
    Рассчет критериев представлен в таблице, все размеры измерялись в байтах.

\begin{tabular}{|c|c|c|c|c|c|}
\hline
№ & \(s_i\) & \( K_4 \) & \( K_5 \) & \( K_6 \) & \( K_7 \) \\
\hline
1 & 1,1,1,1 & 1 & 1 & 0 & 0 \\
\hline
2 & 1,6 & 1 & 6 & 5 & 5\\
\hline
3 & 1,1,2,2,2,3,4,5,5,6 & 1 & 6 & 5& 1 \\
\hline
\end{tabular}

\subsection{Программа}
    Для рассчета тестов была написана простая программа на python
    
\lstinputlisting{cr2.py}
\section{Критерии использующие семантику}
    К данным критериям отнесем критерии использующие дополнительные данные о задаче кроме набора тестов.
\section{Критерии для архивных задач}
    Так как архивные задачи обладают набором посылок с соответствующими вердиктами для таких задач могут быть
    разработаны дополнительные критерии.
\section{Заключение}
\end{document}
